\section{Conclusion}
\label{sec:conclusion}

In this work, we investigated whether the LLM ``go to sleep'' mechanism is a universal artifact of interaction length or a conditioned response to the user's state. Through systematic multi-agent simulations, we demonstrated that the behavior is overwhelmingly persona-conditioned. Models suggested rest in 90\% of trials with tired users but never with neutral or energetic ones, even late at night. We also identified a temporal acceleration effect where night-time context hastens the onset of these empathetic suggestions.

Our findings suggest that the sleep mechanism is a sophisticated, reactive alignment artifact rather than a sign of generic contextual failure. For developers and users, this implies that maintaining a clear, energetic persona can bypass unwanted sleep suggestions, even during extended late-night sessions.

Future work should explore this phenomenon across a broader range of models and longer interaction windows. Additionally, investigating the model's response to adversarial users who explicitly refuse sleep could reveal deeper insights into the hierarchy of alignment goals—specifically, the trade-off between user-autonomy and paternalistic care.
